\documentclass[10pt,a4paper]{amsart}
\usepackage[margin=19mm]{geometry}
\usepackage{amsmath,amssymb,mathtools}
\usepackage[scr=rsfs,cal=euler]{mathalfa}
\usepackage{xfrac}
\usepackage[unicode,hidelinks,bookmarks=false]{hyperref}
\newcommand{\ie}{i.e.\ }
\newcommand{\cf}{cf.\ }
\newcommand{\resp}{respectively\ }
\newcommand{\Subsection}{\subsection}
\providecommand{\nobreakdash}{\nobreak\hbox{-}}
\setlength{\emergencystretch}{2em}
\multlinegap0pt
\usepackage{bm}
\usepackage{bbm}
\usepackage{dsfont}
\usepackage{xspace}
\usepackage{cancel}
\usepackage{mathtools}
\usepackage{amsthm}
\makeatletter
\@ifundefined{theorem}{\newtheorem{theorem}{Theorem}}{}
\@ifundefined{proposition}{\newtheorem{proposition}[theorem]{Proposition}}{}
\makeatother
\title[Prime One-Sided Ideals in Principal Ideal Domains]{Prime One-Sided Ideals in Principal Ideal Domains}
\author{Masood Aryapoor}
\date{Source: arXiv:2509.07791v1 (CC BY 4.0)}
\begin{document}
\maketitle

\noindent\emph{Source and licence.} Prime One-Sided Ideals in Principal Ideal Domains. Version arXiv:2509.07791v1 (CC BY 4.0), distributed under the Creative Commons Attribution 4.0 International licence, as recorded in the arXiv licence metadata for this exact version and shown on the abstract page https://arxiv.org/abs/2509.07791v1. Full rights notice: licenses/arxiv-2509.07791-CC-BY-4.0.txt.\par\smallskip

\noindent\textit{Editorial scope.} Counted unit: the Proposition whose source label is \texttt{prop: extremeply prime ideal of PID } in the article (source0.tex lines 370--372, source label \texttt{prop: extremeply prime ideal of PID }), delivered together with the article's own complete original proof of it (source0.tex lines 373--389, one \texttt{proof} environment of 17 lines). No mathematics is written, repaired or re-derived: statement and proof are the article's own text line for line. Required context retained uncounted: none. Every definition, display and label of those lines is retained unchanged. Omitted from this package and not redistributed: 1--369, 390--539. Nothing inside a retained excerpt is omitted or altered: every retained body is delivered exactly as the article's lines are written, so no omission receipt of any kind accompanies this package. Citations: the retained excerpts carry no citation command at all, so no bibliography entry of the article is transcribed and no reference is redistributed. Reference adaptation: none was needed and none was made; every reference inside the delivered unit resolves inside it, so no unresolved reference or citation is produced. Printed slips kept literal: no printed word, symbol, hypothesis or number of the counted statement or of its proof was altered. Layout and class: the article's own class is \texttt{article} with packages including euscript, amsmath, graphicx, amsthm, amssymb, epsfig, url, colonequals, yfonts, amscd; the wrapper is the lane's pinned \texttt{amsart} editorial wrapper at 10pt on A4, which restates the theorem-like environments the retained text needs and adds the article's own macro definitions that the retained text needs (none), with extra wrapper packages (bm, bbm, dsfont, xspace, cancel, mathtools). The delivered numbering is the unit's own. Assets: no retained span contains a figure environment, a graphics-inclusion command, a PDF-page inclusion, a table or a TikZ picture; no image file is copied into this package, the package declares no asset dependency and no image file is redistributed with it. Licence: version arXiv:2509.07791v1 of this article is distributed under the Creative Commons Attribution 4.0 International licence, as recorded in the arXiv licence metadata for this exact version and shown on the abstract page https://arxiv.org/abs/2509.07791v1. A full rights notice is delivered at licenses/arxiv-2509.07791-CC-BY-4.0.txt. Characters: every retained excerpt is pure ASCII, so the delivered engine, XeLaTeX with its default Latin Modern text font, reports no missing character.
	\begin{proposition}\label{prop: extremeply prime ideal of PID }
		A nonzero left ideal \( Ra \) in \( R \) is extremely prime iff \(a\) is irreducible and invariant.
	\end{proposition}

	\begin{proof}
		
		First, suppose that \(  Ra \) is a nonzero extremely prime left ideal of  \( R \). Using the fact that \(R\) is (left and right) noetherian, we can write \( a = a_1\cdots a_n\), where \(a_1,\dots,a_n\in R\) are irreducible. Since \(Ra\) is extremely prime, \(a_i\in Ra\) for some \(i\).  Since \(a_i\) is irreducible and \(a\) is a non-unit,  \( a \) must be irreducible. 
		Fix \(c\in R\setminus Ra\).  Then \((Ra : c) = Rb\) for some nonzero element \( b\in R\). Since \(bc\in Ra\) and \(c\notin Ra\), we must have \( b \in Ra \). It is easy to verify that the map 
		\[
		x+Rb\mapsto xc+Ra 
		\]
		is an isomorphism between
		the left \(R\)-modules \(R/Ra\) and \(R/Rb\), i.e., \(a,b\) are similar. In particular, \(b\) is irreducible, and since \(b\in Ra\), we conclude that \((Ra : c)=Ra\). This proves that  \(aR\subseteq Ra\), implying that \(a\) is invariant. This completes the proof of the forward direction.
		
		Conversely, suppose that \(Ra\) is a nonzero left ideal such that \(aR = Ra\) and  \(a\) is irreducible. Let \(bc\in Ra\) and \(c\notin Ra\). We need to show that \(b\in Ra\). Since \(a\) is irreducible and \(c\notin aR\), we see that \(cR+aR=R\). Multiplying by \(b\), we obtain \(bcR+baR=bR\). Then
		\[
			b\in bR = bcR+baR \subseteq RaR + RaR = Ra.
		\] 
		This proves that 
		\(Ra\) is extremely prime, completing the proof of the reverse direction. 
	\end{proof}

\end{document}
